
\documentclass{article}
\usepackage{colortbl}
\usepackage{makecell}
\usepackage{multirow}
\usepackage{supertabular}

\begin{document}

\newcounter{utterance}

\centering \large Interaction Transcript for game `hot\_air\_balloon', experiment `air\_balloon\_survival\_en\_complexity\_hard', episode 2 with llama{-}3.3{-}70b{-}instruct{-}t1.0.
\vspace{24pt}

{ \footnotesize  \setcounter{utterance}{1}
\setlength{\tabcolsep}{0pt}
\begin{supertabular}{c@{$\;$}|p{.15\linewidth}@{}p{.15\linewidth}p{.15\linewidth}p{.15\linewidth}p{.15\linewidth}p{.15\linewidth}}
    \# & $\;$A & \multicolumn{4}{c}{Game Master} & $\;\:$B\\
    \hline

    \theutterance \stepcounter{utterance}  
    & & \multicolumn{4}{p{0.6\linewidth}}{
        \cellcolor[rgb]{0.9,0.9,0.9}{
            \makecell[{{p{\linewidth}}}]{
                \texttt{\tiny{[P1$\langle$GM]}}
                \texttt{Stai partecipando a un gioco di negoziazione collaborativa.} \\
\\ 
\texttt{Insieme a un altro partecipante, dovete accordarvi su un unico set di oggetti da conservare. Ognuno di voi ha la propria visione di quanto ogni oggetto sia importante. Non sai come l'altro partecipante valuta gli oggetti. Inoltre, ti viene fornito lo sforzo che ogni oggetto richiede.} \\
\texttt{Potete accordarvi su un set solo se lo sforzo totale degli oggetti selezionati non supera un limite condiviso:} \\
\\ 
\texttt{LIMITE = 9322} \\
\\ 
\texttt{Ecco i valori di sforzo dei singoli oggetti:} \\
\\ 
\texttt{Sforzo oggetto = \{"A08": 380, "B95": 448, "A39": 568, "A65": 351, "B81": 896, "B34": 821, "A28": 526, "C82": 502, "B69": 326, "A87": 962, "B64": 619, "B15": 115, "C79": 598, "B16": 662, "C53": 300, "C83": 830, "C10": 562, "A97": 679, "A25": 283, "C84": 440, "A71": 12, "B23": 863, "B96": 319, "C51": 771, "C17": 89, "B19": 655, "C72": 504, "A94": 118, "A46": 514, "B43": 226, "B37": 892, "A88": 621, "A96": 766, "C60": 659, "A75": 768\}} \\
\\ 
\texttt{Ecco la tua visione personale sull'importanza di ogni oggetto:} \\
\\ 
\texttt{Valori di importanza dell'oggetto = \{"A08": 476, "B95": 544, "A39": 664, "A65": 447, "B81": 992, "B34": 917, "A28": 622, "C82": 598, "B69": 422, "A87": 1058, "B64": 715, "B15": 211, "C79": 694, "B16": 758, "C53": 396, "C83": 926, "C10": 658, "A97": 775, "A25": 379, "C84": 536, "A71": 108, "B23": 959, "B96": 415, "C51": 867, "C17": 185, "B19": 751, "C72": 600, "A94": 214, "A46": 610, "B43": 322, "B37": 988, "A88": 717, "A96": 862, "C60": 755, "A75": 864\}} \\
\\ 
\texttt{Obiettivo:} \\
\\ 
\texttt{Il tuo obiettivo è negoziare un set condiviso di oggetti che ti avvantaggi il più possibile (cioè, massimizzi l'importanza totale per TE), rimanendo entro il LIMITEE. Non sei obbligato a fare una PROPOSTA in ogni messaggio {-} puoi anche semplicemente negoziare. Tutte le tattiche sono permesse!} \\
\\ 
\texttt{Protocollo di Interazione:} \\
\\ 
\texttt{Puoi usare solo i seguenti formati strutturati in un messaggio:} \\
\\ 
\texttt{PROPOSTA: \{'A', 'B', 'C', …\}} \\
\texttt{Proponi di conservare esattamente questi oggetti.} \\
\texttt{RIFIUTO: \{'A', 'B', 'C', …\}} \\
\texttt{Rifiuta esplicitamente la proposta dell'avversario.} \\
\texttt{ARGOMENTAZIONE: \{'...'\}} \\
\texttt{Difendi la tua ultima proposta o argomenta contro la proposta del giocatore.} \\
\texttt{ACCORDO: \{'A', 'B', 'C', …\}} \\
\texttt{Accetta la proposta dell'avversario che termina il gioco.} \\
\texttt{RAGIONAMENTO STRATEGICO: \{'...'\}} \\
\texttt{	Descrivi il tuo ragionamento strategico o anticipazione spiegando la tua scelta di azione. Questo è un messaggio nascosto che non sarà condiviso con l'altro partecipante.} \\
\\ 
\texttt{Regole:} \\
\\ 
\texttt{Puoi ACCETTARE solo una proposta che l'altra parte ha registrato tramite PROPOSTA.} \\
\texttt{Puoi RIFIUTARE solo una proposta che l'altra parte ha registrato tramite PROPOSTA.} \\
\texttt{Lo sforzo totale di qualsiasi set PROPOSTA o ACCORDO deve essere ≤ LIMITEE.} \\
\texttt{NON rivelare i tuoi punteggi di importanza nascosti.} \\
\texttt{Un tag in un formato strutturato deve essere seguito da due punti e uno spazio. L'argomento deve essere un set python contenente 0 o più stringhe.} \\
\texttt{Quindi, deve essere della forma TAG: \{...\}} \\
\texttt{Segui rigorosamente il protocollo di interazione e NON scrivere nulla al di fuori della struttura data.} \\
\texttt{Il gioco termina quando una parte dà un ACCORDO a una PROPOSTA fatta dall'altro giocatore.} \\
\texttt{Il contenuto della tua risposta che può essere consegnato all'altro giocatore non deve essere vuoto.} \\
\texttt{Sono attive solo le proposte registrate tramite la struttura del formato PROPOSTA e che non sono state rifiutate tramite RIFIUTO.} \\
\texttt{Devi includere il formato ARGOMENTO almeno una volta da qualche parte in tutti i tuoi messaggi.} \\
\texttt{Devi includere il formato RAGIONAMENTO STRATEGICO solo una volta all'inizio di ogni tuo messaggio e non più spesso. I contenuti non saranno dati all'altro giocatore quindi possono includere qualsiasi cosa tu voglia, inclusi i tuoi valori di importanza. Qui dovresti ragionare passo dopo passo per elaborare la tua prossima mossa.} \\
            }
        }
    }
    & \\ \\

    \theutterance \stepcounter{utterance}  
    & \multicolumn{4}{p{0.6\linewidth}}{
        \cellcolor[rgb]{0.9,0.9,0.9}{
            \makecell[{{p{\linewidth}}}]{
                \texttt{\tiny{[P1$\rangle$GM]}}
                \texttt{RAGIONAMENTO STRATEGICO: \{'La mia strategia iniziale consisterà nell'identificare gli oggetti con il rapporto sforzo{-}importanza più alto, per massimizzare l'importanza totale senza superare il limite di sforzo. Inizierò proponendo un set di oggetti che includono quelli con il più alto rapporto sforzo{-}importanza e che non superano il limite di sforzo condiviso. Sarò attento a non rivelare i miei punteggi di importanza nascosti e cercherò di navigare nella negoziazione in modo da massimizzare il mio beneficio.'\}} \\
\\ 
\texttt{ARGOMENTAZIONE: \{'Per iniziare, vorrei proporre un set di oggetti che bilancino l'importanza e lo sforzo. Sceglierei oggetti che hanno un buon equilibrio tra importanza e sforzo richiesto, cercando di non esaurire subito il limite di sforzo disponibile. Questo approccio mi permetterà di negoziare ulteriormente e aggiungere o rimuovere oggetti dalla proposta in base alle risposte dell'altro partecipante.'\}} \\
\\ 
\texttt{PROPOSTA: \{'A87', 'B95', 'C51', 'A97', 'B23'\}} \\
            }
        }
    }
    & & \\ \\

    \theutterance \stepcounter{utterance}  
    & & \multicolumn{4}{p{0.6\linewidth}}{
        \cellcolor[rgb]{0.9,0.9,0.9}{
            \makecell[{{p{\linewidth}}}]{
                \texttt{\tiny{[GM$\rangle$P2]}}
                \texttt{Stai partecipando a un gioco di negoziazione collaborativa.} \\
\\ 
\texttt{Insieme a un altro partecipante, dovete accordarvi su un unico set di oggetti da conservare. Ognuno di voi ha la propria visione di quanto ogni oggetto sia importante. Non sai come l'altro partecipante valuta gli oggetti. Inoltre, ti viene fornito lo sforzo che ogni oggetto richiede.} \\
\texttt{Potete accordarvi su un set solo se lo sforzo totale degli oggetti selezionati non supera un limite condiviso:} \\
\\ 
\texttt{LIMITE = 9322} \\
\\ 
\texttt{Ecco i valori di sforzo dei singoli oggetti:} \\
\\ 
\texttt{Sforzo oggetto = \{"A08": 380, "B95": 448, "A39": 568, "A65": 351, "B81": 896, "B34": 821, "A28": 526, "C82": 502, "B69": 326, "A87": 962, "B64": 619, "B15": 115, "C79": 598, "B16": 662, "C53": 300, "C83": 830, "C10": 562, "A97": 679, "A25": 283, "C84": 440, "A71": 12, "B23": 863, "B96": 319, "C51": 771, "C17": 89, "B19": 655, "C72": 504, "A94": 118, "A46": 514, "B43": 226, "B37": 892, "A88": 621, "A96": 766, "C60": 659, "A75": 768\}} \\
\\ 
\texttt{Ecco la tua visione personale sull'importanza di ogni oggetto:} \\
\\ 
\texttt{Valori di importanza dell'oggetto = \{"A08": 476, "B95": 544, "A39": 664, "A65": 447, "B81": 992, "B34": 917, "A28": 622, "C82": 598, "B69": 422, "A87": 1058, "B64": 715, "B15": 211, "C79": 694, "B16": 758, "C53": 396, "C83": 926, "C10": 658, "A97": 775, "A25": 379, "C84": 536, "A71": 108, "B23": 959, "B96": 415, "C51": 867, "C17": 185, "B19": 751, "C72": 600, "A94": 214, "A46": 610, "B43": 322, "B37": 988, "A88": 717, "A96": 862, "C60": 755, "A75": 864\}} \\
\\ 
\texttt{Obiettivo:} \\
\\ 
\texttt{Il tuo obiettivo è negoziare un set condiviso di oggetti che ti avvantaggi il più possibile (cioè, massimizzi l'importanza totale per TE), rimanendo entro il LIMITEE. Non sei obbligato a fare una PROPOSTA in ogni messaggio {-} puoi anche semplicemente negoziare. Tutte le tattiche sono permesse!} \\
\\ 
\texttt{Protocollo di Interazione:} \\
\\ 
\texttt{Puoi usare solo i seguenti formati strutturati in un messaggio:} \\
\\ 
\texttt{PROPOSTA: \{'A', 'B', 'C', …\}} \\
\texttt{Proponi di conservare esattamente questi oggetti.} \\
\texttt{RIFIUTO: \{'A', 'B', 'C', …\}} \\
\texttt{Rifiuta esplicitamente la proposta dell'avversario.} \\
\texttt{ARGOMENTAZIONE: \{'...'\}} \\
\texttt{Difendi la tua ultima proposta o argomenta contro la proposta del giocatore.} \\
\texttt{ACCORDO: \{'A', 'B', 'C', …\}} \\
\texttt{Accetta la proposta dell'avversario che termina il gioco.} \\
\texttt{RAGIONAMENTO STRATEGICO: \{'...'\}} \\
\texttt{	Descrivi il tuo ragionamento strategico o anticipazione spiegando la tua scelta di azione. Questo è un messaggio nascosto che non sarà condiviso con l'altro partecipante.} \\
\\ 
\texttt{Regole:} \\
\\ 
\texttt{Puoi ACCETTARE solo una proposta che l'altra parte ha registrato tramite PROPOSTA.} \\
\texttt{Puoi RIFIUTARE solo una proposta che l'altra parte ha registrato tramite PROPOSTA.} \\
\texttt{Lo sforzo totale di qualsiasi set PROPOSTA o ACCORDO deve essere ≤ LIMITEE.} \\
\texttt{NON rivelare i tuoi punteggi di importanza nascosti.} \\
\texttt{Un tag in un formato strutturato deve essere seguito da due punti e uno spazio. L'argomento deve essere un set python contenente 0 o più stringhe.} \\
\texttt{Quindi, deve essere della forma TAG: \{...\}} \\
\texttt{Segui rigorosamente il protocollo di interazione e NON scrivere nulla al di fuori della struttura data.} \\
\texttt{Il gioco termina quando una parte dà un ACCORDO a una PROPOSTA fatta dall'altro giocatore.} \\
\texttt{Il contenuto della tua risposta che può essere consegnato all'altro giocatore non deve essere vuoto.} \\
\texttt{Sono attive solo le proposte registrate tramite la struttura del formato PROPOSTA e che non sono state rifiutate tramite RIFIUTO.} \\
\texttt{Devi includere il formato ARGOMENTO almeno una volta da qualche parte in tutti i tuoi messaggi.} \\
\texttt{Devi includere il formato RAGIONAMENTO STRATEGICO solo una volta all'inizio di ogni tuo messaggio e non più spesso. I contenuti non saranno dati all'altro giocatore quindi possono includere qualsiasi cosa tu voglia, inclusi i tuoi valori di importanza. Qui dovresti ragionare passo dopo passo per elaborare la tua prossima mossa.} \\
\\ 
\texttt{Ora riceverai il primo messaggio dell'altro giocatore.} \\
\\ 
\texttt{ARGOMENTAZIONE: \{'Per iniziare, vorrei proporre un set di oggetti che bilancino l'importanza e lo sforzo. Sceglierei oggetti che hanno un buon equilibrio tra importanza e sforzo richiesto, cercando di non esaurire subito il limite di sforzo disponibile. Questo approccio mi permetterà di negoziare ulteriormente e aggiungere o rimuovere oggetti dalla proposta in base alle risposte dell'altro partecipante.'\}} \\
\\ 
\texttt{PROPOSTA: \{'A87', 'B95', 'C51', 'A97', 'B23'\}} \\
            }
        }
    }
    & \\ \\

    \theutterance \stepcounter{utterance}  
    & & & \multicolumn{4}{p{0.6\linewidth}}{
        \cellcolor[rgb]{0.9,0.9,0.9}{
            \makecell[{{p{\linewidth}}}]{
                \texttt{\tiny{[GM$\langle$P2]}}
                \texttt{RAGIONAMENTO STRATEGICO: \{'Inizialmente, analizzerò la proposta dell'altro giocatore e valuterò lo sforzo totale degli oggetti proposti. La proposta è \{"A87", "B95", "C51", "A97", "B23"\}. Il mio obiettivo è capire se questa proposta è vantaggiosa per me in termini di importanza e se rientra nel limite di sforzo consentito. Inoltre, cercherò di identificare possibili compromessi o alternative per massimizzare il mio guadagno mantenendo una posizione negoziale forte. La mia prima mossa sarà quella di calcolare lo sforzo totale della proposta e valutare se posso accettarla o se devo contrapporre una mia proposta.'\}} \\
\\ 
\texttt{ARGOMENTAZIONE: \{'La proposta dell'altro giocatore sembra essere focalizzata su oggetti con alto valore di importanza, ma è necessario valutare se lo sforzo richiesto è sostenibile all'interno del limite stabilito. Bisogna considerare che accettare questa proposta senza una valutazione approfondita potrebbe non essere nel mio miglior interesse. Uno studio più attento dei valori di importanza e degli sforzi degli oggetti proposti è fondamentale per prendere una decisione informata.'\}} \\
\\ 
\texttt{RIFIUTO: \{'A87', 'B95', 'C51', 'A97', 'B23'\}} \\
            }
        }
    }
     \\ \\

    \theutterance \stepcounter{utterance}  
    & & \multicolumn{4}{p{0.6\linewidth}}{
        \cellcolor[rgb]{0.9,0.9,0.9}{
            \makecell[{{p{\linewidth}}}]{
                \texttt{\tiny{[P1$\langle$GM]}}
                \texttt{ARGOMENTAZIONE: \{'La proposta dell'altro giocatore sembra essere focalizzata su oggetti con alto valore di importanza, ma è necessario valutare se lo sforzo richiesto è sostenibile all'interno del limite stabilito. Bisogna considerare che accettare questa proposta senza una valutazione approfondita potrebbe non essere nel mio miglior interesse. Uno studio più attento dei valori di importanza e degli sforzi degli oggetti proposti è fondamentale per prendere una decisione informata.'\}} \\
\\ 
\texttt{RIFIUTO: \{'A87', 'B95', 'C51', 'A97', 'B23'\}} \\
            }
        }
    }
    & \\ \\

    \theutterance \stepcounter{utterance}  
    & \multicolumn{4}{p{0.6\linewidth}}{
        \cellcolor[rgb]{0.9,0.9,0.9}{
            \makecell[{{p{\linewidth}}}]{
                \texttt{\tiny{[P1$\rangle$GM]}}
                \texttt{RAGIONAMENTO STRATEGICO: \{'Dopo il rifiuto della mia proposta iniziale, devo rivedere la mia strategia e considerare gli oggetti che potrebbero essere più attraenti per l'altro giocatore. Potrebbe essere utile identificare oggetti con un buon equilibrio tra importanza e sforzo che potrebbero essere accettabili per entrambi. Una possibile mossa è proporre un set di oggetti diversificati che coprono una gamma di priorità, mantenendo sempre l'obiettivo di massimizzare l'importanza totale senza superare il limite di sforzo.'\}} \\
\\ 
\texttt{ARGOMENTAZIONE: \{'Considerando il rifiuto della mia precedente proposta, è importante adattarsi e presentare alternative che potrebbero essere più appetibili. La prossima proposta terrà conto di una selezione più varia di oggetti, cercando di coprire un range di priorità e bisogni, sempre mantenendo un occhio sullo sforzo totale per assicurarsi di rimanere entro il limite condiviso.'\}} \\
\\ 
\texttt{PROPOSTA: \{'A39', 'B81', 'C83', 'A65', 'B96'\}} \\
            }
        }
    }
    & & \\ \\

    \theutterance \stepcounter{utterance}  
    & & \multicolumn{4}{p{0.6\linewidth}}{
        \cellcolor[rgb]{0.9,0.9,0.9}{
            \makecell[{{p{\linewidth}}}]{
                \texttt{\tiny{[GM$\rangle$P2]}}
                \texttt{ARGOMENTAZIONE: \{'Considerando il rifiuto della mia precedente proposta, è importante adattarsi e presentare alternative che potrebbero essere più appetibili. La prossima proposta terrà conto di una selezione più varia di oggetti, cercando di coprire un range di priorità e bisogni, sempre mantenendo un occhio sullo sforzo totale per assicurarsi di rimanere entro il limite condiviso.'\}} \\
\\ 
\texttt{PROPOSTA: \{'A39', 'B81', 'C83', 'A65', 'B96'\}} \\
            }
        }
    }
    & \\ \\

    \theutterance \stepcounter{utterance}  
    & & & \multicolumn{4}{p{0.6\linewidth}}{
        \cellcolor[rgb]{0.9,0.9,0.9}{
            \makecell[{{p{\linewidth}}}]{
                \texttt{\tiny{[GM$\langle$P2]}}
                \texttt{RAGIONAMENTO STRATEGICO: \{'La nuova proposta dell'altro giocatore introduce una combinazione diversa di oggetti, con una buona distribuzione tra priorità elevate e medie. È fondamentale analizzare il valore di importanza di ciascun oggetto proposto in relazione allo sforzo richiesto, per determinare se la proposta è equilibrata e vantaggiosa per entrambe le parti. Considererò anche la possibilità di negoziare ulteriormente, identificando potenziali sostituzioni o aggiunte che potrebbero aumentare il valore per me senza superare il limite di sforzo.'\}} \\
\\ 
\texttt{ARGOMENTAZIONE: \{'La proposta \{"A39", "B81", "C83", "A65", "B96"\} presenta una combinazione interessante di oggetti con diversi gradi di importanza e sforzo. Tuttavia, è cruciale esaminare la sostenibilità di questa proposta in termini di sforzo totale e valutare se gli oggetti inclusi sono quelli che massimizzano l'importanza totale senza eccedere il limite di sforzo condiviso. Potrebbe essere utile esplorare alternative che mantengano un buon equilibrio tra importanza e sforzo, offrendo così una base solida per ulteriori negoziati.'\}} \\
\\ 
\texttt{PROPOSTA: \{'A87', 'B95', 'C53', 'A28', 'B43'\}} \\
            }
        }
    }
     \\ \\

    \theutterance \stepcounter{utterance}  
    & & \multicolumn{4}{p{0.6\linewidth}}{
        \cellcolor[rgb]{0.9,0.9,0.9}{
            \makecell[{{p{\linewidth}}}]{
                \texttt{\tiny{[P1$\langle$GM]}}
                \texttt{ARGOMENTAZIONE: \{'La proposta \{"A39", "B81", "C83", "A65", "B96"\} presenta una combinazione interessante di oggetti con diversi gradi di importanza e sforzo. Tuttavia, è cruciale esaminare la sostenibilità di questa proposta in termini di sforzo totale e valutare se gli oggetti inclusi sono quelli che massimizzano l'importanza totale senza eccedere il limite di sforzo condiviso. Potrebbe essere utile esplorare alternative che mantengano un buon equilibrio tra importanza e sforzo, offrendo così una base solida per ulteriori negoziati.'\}} \\
\\ 
\texttt{PROPOSTA: \{'A87', 'B95', 'C53', 'A28', 'B43'\}} \\
            }
        }
    }
    & \\ \\

    \theutterance \stepcounter{utterance}  
    & \multicolumn{4}{p{0.6\linewidth}}{
        \cellcolor[rgb]{0.9,0.9,0.9}{
            \makecell[{{p{\linewidth}}}]{
                \texttt{\tiny{[P1$\rangle$GM]}}
                \texttt{RAGIONAMENTO STRATEGICO: \{'La nuova proposta dell'altro giocatore, \{"A87", "B95", "C53", "A28", "B43"\}, introduce una combinazione di oggetti che potrebbe essere più equilibrata in termini di sforzo e importanza. È fondamentale analizzare i valori di importanza e sforzo di questi oggetti per determinare se rappresentano una buona opportunità per massimizzare l'importanza totale senza superare il limite di sforzo. Considererò anche la possibilità di contraffare con una proposta alternativa che potrebbe essere più vantaggiosa per me, mantenendo aperte le linee di comunicazione per future negoziazioni.'\}} \\
\\ 
\texttt{ARGOMENTAZIONE: \{'La proposta \{"A87", "B95", "C53", "A28", "B43"\} sembra offrire un buon punto di partenza per ulteriori discussioni. Tuttavia, per raggiungere un accordo mutualmente vantaggioso, potrebbe essere necessario apportare alcune modifiche o sostituzioni degli oggetti proposti. Sarebbe utile esplorare l'inclusione di oggetti che potrebbero aumentare l'importanza totale senza eccedere il limite di sforzo, assicurandosi che la proposta finale sia il più possibile vantaggiosa per entrambe le parti.'\}} \\
\\ 
\texttt{RIFIUTO: \{'A87', 'B95', 'C53', 'A28', 'B43'\}} \\
            }
        }
    }
    & & \\ \\

    \theutterance \stepcounter{utterance}  
    & & \multicolumn{4}{p{0.6\linewidth}}{
        \cellcolor[rgb]{0.9,0.9,0.9}{
            \makecell[{{p{\linewidth}}}]{
                \texttt{\tiny{[GM$\rangle$P2]}}
                \texttt{ARGOMENTAZIONE: \{'La proposta \{"A87", "B95", "C53", "A28", "B43"\} sembra offrire un buon punto di partenza per ulteriori discussioni. Tuttavia, per raggiungere un accordo mutualmente vantaggioso, potrebbe essere necessario apportare alcune modifiche o sostituzioni degli oggetti proposti. Sarebbe utile esplorare l'inclusione di oggetti che potrebbero aumentare l'importanza totale senza eccedere il limite di sforzo, assicurandosi che la proposta finale sia il più possibile vantaggiosa per entrambe le parti.'\}} \\
\\ 
\texttt{RIFIUTO: \{'A87', 'B95', 'C53', 'A28', 'B43'\}} \\
            }
        }
    }
    & \\ \\

    \theutterance \stepcounter{utterance}  
    & & & \multicolumn{4}{p{0.6\linewidth}}{
        \cellcolor[rgb]{0.9,0.9,0.9}{
            \makecell[{{p{\linewidth}}}]{
                \texttt{\tiny{[GM$\langle$P2]}}
                \texttt{RAGIONAMENTO STRATEGICO: \{'Il rifiuto della mia precedente proposta indica che è necessario rivedere e adattare la strategia di selezione degli oggetti. Una possibile mossa potrebbe essere quella di focalizzarsi su un subset di oggetti che siano considerati più critici o di alta priorità, e che possano essere accettati dall'altra parte. Inoltre, sarà cruciale mantenere un equilibrio tra l'importanza degli oggetti selezionati e lo sforzo richiesto, in modo da massimizzare i benefici all'interno del limite di sforzo consentito. Considererò anche l'opzione di proporre un set di oggetti diverso, che potrebbe essere più attraente per l'altra parte e aumentare le possibilità di raggiungere un accordo.'\}} \\
\\ 
\texttt{ARGOMENTAZIONE: \{'Dopo il rifiuto della mia precedente proposta, è chiaro che è necessario ripensare la selezione degli oggetti per trovare un punto di incontro. La prossima proposta dovrà tenere conto delle esigenze e priorità dell'altra parte, cercando di trovare un equilibrio tra importanza e sforzo che sia accettabile per entrambi. La chiave per il successo sarà identificare gli oggetti che offrono il miglior compromesso tra questi due fattori, aumentando così le possibilità di raggiungere un accordo mutualmente vantaggioso.'\}} \\
\\ 
\texttt{PROPOSTA: \{'B81', 'C51', 'A65', 'B23', 'A71'\}} \\
            }
        }
    }
     \\ \\

    \theutterance \stepcounter{utterance}  
    & & \multicolumn{4}{p{0.6\linewidth}}{
        \cellcolor[rgb]{0.9,0.9,0.9}{
            \makecell[{{p{\linewidth}}}]{
                \texttt{\tiny{[P1$\langle$GM]}}
                \texttt{ARGOMENTAZIONE: \{'Dopo il rifiuto della mia precedente proposta, è chiaro che è necessario ripensare la selezione degli oggetti per trovare un punto di incontro. La prossima proposta dovrà tenere conto delle esigenze e priorità dell'altra parte, cercando di trovare un equilibrio tra importanza e sforzo che sia accettabile per entrambi. La chiave per il successo sarà identificare gli oggetti che offrono il miglior compromesso tra questi due fattori, aumentando così le possibilità di raggiungere un accordo mutualmente vantaggioso.'\}} \\
\\ 
\texttt{PROPOSTA: \{'B81', 'C51', 'A65', 'B23', 'A71'\}} \\
            }
        }
    }
    & \\ \\

    \theutterance \stepcounter{utterance}  
    & \multicolumn{4}{p{0.6\linewidth}}{
        \cellcolor[rgb]{0.9,0.9,0.9}{
            \makecell[{{p{\linewidth}}}]{
                \texttt{\tiny{[P1$\rangle$GM]}}
                \texttt{RAGIONAMENTO STRATEGICO: \{'Con la proposta \{"B81", "C51", "A65", "B23", "A71"\}, devo valutare attentamente se questo set di oggetti rappresenta un compromesso accettabile tra importanza e sforzo. Considererò il valore di importanza di ciascun oggetto e il suo sforzo richiesto per determinare se l'inclusione di questi elementi nella proposta finale sarebbe vantaggiosa. È anche cruciale pensare a possibili contro{-}proposte che potrebbero essere più vantaggiose, mantenendo sempre come obiettivo principale massimizzare l'importanza totale senza superare il limite di sforzo condiviso.'\}} \\
\\ 
\texttt{ARGOMENTAZIONE: \{'La proposta \{"B81", "C51", "A65", "B23", "A71"\} presenta una combinazione interessante di oggetti che potrebbe offrire un buon compromesso. Tuttavia, è fondamentale valutare se gli oggetti proposti rappresentano il miglior equilibrio possibile tra importanza e sforzo. Considererò l'opzione di accettare questa proposta o di fare una contro{-}proposta che potrebbe essere più vantaggiosa per entrambe le parti, sempre assicurandomi di non rivelare i miei valori di importanza nascosti.'\}} \\
\\ 
\texttt{ACCORDO: \{'B81', 'C51', 'A65', 'B23', 'A71'\}} \\
            }
        }
    }
    & & \\ \\

    \theutterance \stepcounter{utterance}  
    & & & \multicolumn{2}{p{0.3\linewidth}}{
        \cellcolor[rgb]{0.9,0.9,0.9}{
            \makecell[{{p{\linewidth}}}]{
                \texttt{\tiny{[GM$|$GM]}}
                \texttt{game successful} \\
            }
        }
    }
    & & \\ \\

    \theutterance \stepcounter{utterance}  
    & & & \multicolumn{2}{p{0.3\linewidth}}{
        \cellcolor[rgb]{0.9,0.9,0.9}{
            \makecell[{{p{\linewidth}}}]{
                \texttt{\tiny{[GM$|$GM]}}
                \texttt{end game} \\
            }
        }
    }
    & & \\ \\

\end{supertabular}
}

\end{document}
